\section{Banco de questões — Noções de Direitos Humanos}

\subsection*{N1 — Fundamentos em nível de concurso}
\begin{questao}{\nivel{N1} 09.01 — Cebraspe (C/E)}A distinção entre direitos humanos e direitos fundamentais é absoluta e impede que tratados internacionais influenciem a interpretação constitucional interna.\end{questao}
\begin{questao}{\nivel{N1} 09.02 — Cebraspe (C/E)}A classificação dos direitos em dimensões não significa substituição dos direitos de liberdade pelos direitos sociais ou coletivos.\end{questao}
\begin{questao}{\nivel{N1} 09.03 — FGV}A Declaração Universal dos Direitos Humanos foi proclamada:\begin{enumerate}[label=\Alph*)]\item em 1948 pela Assembleia Geral da ONU.\item em 1988 pelo Congresso Nacional brasileiro.\item em 2015 junto com a Agenda 2030.\item em 1969 pela OEA.\item em 1945 pelo STF.\end{enumerate}\end{questao}
\begin{questao}{\nivel{N1} 09.04 — Cebraspe (C/E)}A DUDH contempla direitos civis e políticos e também direitos econômicos, sociais e culturais.\end{questao}
\begin{questao}{\nivel{N1} 09.05 — Cebraspe (C/E)}A Agenda 2030 organiza-se em 17 Objetivos de Desenvolvimento Sustentável e 169 metas.\end{questao}
\begin{questao}{\nivel{N1} 09.06 — FCC}Segundo a LBI, pessoa com deficiência é definida a partir:\begin{enumerate}[label=\Alph*)]\item apenas do diagnóstico médico.\item da interação entre impedimento de longo prazo e barreiras que podem obstruir participação plena e efetiva.\item exclusivamente da incapacidade civil.\item apenas da existência de benefício previdenciário.\item somente da idade.\end{enumerate}\end{questao}
\begin{questao}{\nivel{N1} 09.07 — Cebraspe (C/E)}A deficiência, por si só, não afeta a plena capacidade civil da pessoa.\end{questao}
\begin{questao}{\nivel{N1} 09.08 — Cebraspe (C/E)}A recusa de adaptação razoável pode configurar discriminação em razão da deficiência.\end{questao}
\begin{questao}{\nivel{N1} 09.09 — Vunesp}A Lei nº 10.098/2000 trata principalmente de:\begin{enumerate}[label=\Alph*)]\item acessibilidade.\item improbidade administrativa.\item licitações.\item orçamento secreto.\item processo penal.\end{enumerate}\end{questao}
\begin{questao}{\nivel{N1} 09.10 — Cebraspe (C/E)}O Estatuto da Igualdade Racial reconhece ações afirmativas como medidas voltadas à correção de desigualdades raciais e promoção de igualdade de oportunidades.\end{questao}

\subsection*{N2 — Aplicação}
\begin{questao}{\nivel{N2} 09.11 — FGV}Um sistema público é fisicamente acessível em seu posto presencial, mas o portal digital não pode ser usado por leitor de tela. Nesse caso:\begin{enumerate}[label=\Alph*)]\item não há problema porque acessibilidade se limita ao espaço físico.\item existe barreira de comunicação/informação e tecnologia.\item acessibilidade digital é facultativa para o poder público.\item o problema só existiria se o prédio não tivesse elevador.\item a LBI não alcança tecnologia.\end{enumerate}\end{questao}
\begin{questao}{\nivel{N2} 09.12 — Cebraspe (C/E)}A avaliação biopsicossocial, quando necessária, considera fatores que vão além das funções corporais, incluindo ambiente, atividades e participação.\end{questao}
\begin{questao}{\nivel{N2} 09.13 — Cebraspe (C/E)}Tratar todas as pessoas de modo rigorosamente idêntico é sempre suficiente para realizar o princípio da igualdade material.\end{questao}
\begin{questao}{\nivel{N2} 09.14 — FCC}Qual alternativa melhor exemplifica adaptação razoável?\begin{enumerate}[label=\Alph*)]\item negar acesso por inexistência de padrão único.\item realizar ajuste necessário e adequado, sem ônus desproporcional ou indevido, para assegurar exercício de direito.\item retirar autonomia da pessoa com deficiência.\item exigir que o usuário se adapte ao serviço inacessível.\item substituir acessibilidade por atendimento segregado obrigatório.\end{enumerate}\end{questao}
\begin{questao}{\nivel{N2} 09.15 — Cebraspe (C/E)}A existência de deficiência autoriza presumir incapacidade para casar, exercer direitos reprodutivos ou decidir sobre a própria vida.\end{questao}
\begin{questao}{\nivel{N2} 09.16 — FGV}A Lei nº 10.048/2000, em sua redação atual, inclui entre os beneficiários de atendimento prioritário:\begin{enumerate}[label=\Alph*)]\item apenas pessoas com deficiência física.\item pessoas com deficiência, TEA, idosos de 60 anos ou mais, gestantes, lactantes e outros grupos previstos em lei.\item somente servidores públicos.\item exclusivamente gestantes e idosos de 65 anos ou mais.\item apenas pessoas com laudo judicial.\end{enumerate}\end{questao}
\begin{questao}{\nivel{N2} 09.17 — Cebraspe (C/E)}Doadores de sangue possuem prioridade legal exatamente na mesma posição dos demais beneficiários e sem necessidade de qualquer comprovação.\end{questao}
\begin{questao}{\nivel{N2} 09.18 — Cebraspe (C/E)}A Agenda 2030 é universal e se aplica a países desenvolvidos e em desenvolvimento.\end{questao}
\begin{questao}{\nivel{N2} 09.19 — FGV}O ODS mais diretamente relacionado a instituições eficazes, responsáveis e inclusivas e acesso à justiça é o:\begin{enumerate}[label=\Alph*)]\item ODS 2.\item ODS 5.\item ODS 10.\item ODS 16.\item ODS 17 apenas.\end{enumerate}\end{questao}
\begin{questao}{\nivel{N2} 09.20 — Cebraspe (C/E)}Os ODS devem ser compreendidos como integrados e indivisíveis, e não como objetivos totalmente independentes.\end{questao}
\begin{questao}{\nivel{N2} 09.21 — FCC}O Estatuto da Igualdade Racial define população negra como:\begin{enumerate}[label=\Alph*)]\item apenas pessoas autodeclaradas pretas.\item conjunto de pessoas autodeclaradas pretas e pardas, segundo quesito de cor ou raça utilizado pelo IBGE, ou autodefinição análoga.\item apenas pessoas com certidão específica.\item qualquer estrangeiro residente no Brasil.\item somente comunidades quilombolas.\end{enumerate}\end{questao}
\begin{questao}{\nivel{N2} 09.22 — Cebraspe (C/E)}Ações afirmativas podem ser compatíveis com igualdade porque buscam corrigir desigualdades concretas e ampliar oportunidades.\end{questao}
\begin{questao}{\nivel{N2} 09.23 — Cebraspe (C/E)}Barreira atitudinal pode existir mesmo em ambiente arquitetonicamente acessível.\end{questao}
\begin{questao}{\nivel{N2} 09.24 — FGV}Um órgão publica vídeos oficiais sem legendas ou recurso equivalente e mantém formulários incompatíveis com tecnologia assistiva. O problema mais adequado é classificado como:\begin{enumerate}[label=\Alph*)]\item barreira de comunicação/informação e tecnológica.\item ausência de competência tributária.\item falha exclusivamente urbanística.\item problema sem relação com direitos humanos.\item exercício legítimo de discriminação.\end{enumerate}\end{questao}
\begin{questao}{\nivel{N2} 09.25 — Cebraspe (C/E)}A Agenda 2030 busca integrar desenvolvimento econômico, inclusão social e proteção ambiental.\end{questao}

\subsection*{N3 — Nível de prova}
\begin{questao}{\nivel{N3} 09.26 — Cebraspe (C/E)}A universalidade dos direitos humanos significa que todo direito é absoluto e jamais admite ponderação ou restrição constitucionalmente legítima.\end{questao}
\begin{questao}{\nivel{N3} 09.27 — FGV}Um edital exige de todo candidato com deficiência curatela judicial para reconhecer sua capacidade de inscrição. A exigência é problemática porque:\begin{enumerate}[label=\Alph*)]\item deficiência não implica incapacidade civil automática.\item toda pessoa com deficiência é incapaz.\item curatela é requisito geral de cidadania.\item a LBI eliminou a personalidade jurídica.\item concursos não se submetem à igualdade.\end{enumerate}\end{questao}
\begin{questao}{\nivel{N3} 09.28 — Cebraspe (C/E)}A curatela, na lógica protetiva contemporânea, não deve operar como mecanismo genérico de supressão de direitos existenciais da pessoa com deficiência.\end{questao}
\begin{questao}{\nivel{N3} 09.29 — Cebraspe (C/E)}Acessibilidade digital pode ser auditada mediante testes de navegação por teclado, compatibilidade com leitores de tela, contraste, rótulos e estrutura semântica.\end{questao}
\begin{questao}{\nivel{N3} 09.30 — FCC}Uma política pública oferece o mesmo canal exclusivamente visual a todos os cidadãos. Pessoas cegas não conseguem utilizá-lo. Sob igualdade material, a solução adequada é:\begin{enumerate}[label=\Alph*)]\item manter o canal porque igualdade exige identidade formal.\item remover barreiras e oferecer recursos acessíveis equivalentes.\item excluir usuários com deficiência.\item condicionar acesso a representante.\item impedir uso de tecnologia assistiva.\end{enumerate}\end{questao}
\begin{questao}{\nivel{N3} 09.31 — Cebraspe (C/E)}A recusa injustificada de tecnologia assistiva necessária pode integrar o conceito de discriminação em razão da deficiência.\end{questao}
\begin{questao}{\nivel{N3} 09.32 — FGV}Sobre a DUDH, assinale a opção correta:\begin{enumerate}[label=\Alph*)]\item restringe-se a direitos civis.\item admite escravidão em hipóteses nacionais.\item consagra dignidade, igualdade e conjunto amplo de direitos.\item foi aprovada como emenda à Constituição brasileira.\item não aborda educação ou trabalho.\end{enumerate}\end{questao}
\begin{questao}{\nivel{N3} 09.33 — Cebraspe (C/E)}A liberdade de expressão prevista na DUDH integra conjunto de direitos que também inclui liberdade de pensamento, consciência e religião.\end{questao}
\begin{questao}{\nivel{N3} 09.34 — Cebraspe (C/E)}O ODS 16 possui conexão direta com transparência, acesso à informação e instituições responsáveis, o que o torna relevante para controle externo.\end{questao}
\begin{questao}{\nivel{N3} 09.35 — FGV}Uma auditoria avalia se programa reduziu desigualdade de acesso entre grupos raciais, mas o banco de dados não contém variável de raça/cor. A principal limitação é:\begin{enumerate}[label=\Alph*)]\item impossibilidade de desagregar resultados e medir desigualdade racial.\item excesso de informação pessoal necessariamente ilícita.\item ausência de relação com políticas de igualdade.\item impossibilidade de qualquer avaliação de desempenho.\item necessidade de excluir todo dado demográfico.\end{enumerate}\end{questao}
\begin{questao}{\nivel{N3} 09.36 — Cebraspe (C/E)}Igualdade racial exige apenas proibição formal de discriminação, sendo incompatível com políticas públicas específicas voltadas a corrigir desigualdades persistentes.\end{questao}
\begin{questao}{\nivel{N3} 09.37 — Cebraspe (C/E)}O Estatuto da Igualdade Racial alcança direitos em áreas como educação, cultura, trabalho, liberdade religiosa e acesso à justiça.\end{questao}
\begin{questao}{\nivel{N3} 09.38 — FGV}Uma repartição possui fila prioritária, mas impede acompanhante indispensável de auxiliar pessoa com deficiência durante atendimento. À luz da proteção vigente, a conduta deve ser analisada considerando:\begin{enumerate}[label=\Alph*)]\item autonomia, apoio e efetividade do atendimento prioritário.\item apenas ordem de chegada.\item exclusão obrigatória do acompanhante.\item inexistência de prioridade legal.\item impossibilidade de ajuste do atendimento.\end{enumerate}\end{questao}
\begin{questao}{\nivel{N3} 09.39 — Cebraspe (C/E)}O cordão de girassóis é símbolo de uso opcional para identificação de deficiências ocultas; sua ausência não elimina direitos legalmente assegurados.\end{questao}
\begin{questao}{\nivel{N3} 09.40 — Cebraspe (C/E)}A legislação atual passou a contemplar necessidades complexas de comunicação e soluções de comunicação aumentativa e alternativa em determinados contextos.\end{questao}

\subsection*{N4 — Avançado e controle}
\begin{questao}{\nivel{N4} 09.41 — FGV}Um órgão compra sistema de R\$ 20 milhões. O termo de referência exige disponibilidade e segurança, mas nada sobre acessibilidade. Após implantação, usuários cegos não conseguem usar o serviço. O achado mais completo relaciona:\begin{enumerate}[label=\Alph*)]\item falha de requisitos, barreira digital e risco de discriminação/exclusão no serviço público.\item apenas defeito cosmético.\item exclusivamente indisponibilidade.\item problema sem impacto jurídico.\item excesso de segurança.\end{enumerate}\end{questao}
\begin{questao}{\nivel{N4} 09.42 — Cebraspe (C/E)}Em auditoria de política inclusiva, avaliar somente volume total de beneficiários pode ocultar desigualdades, razão pela qual dados desagregados e indicadores de acesso são relevantes.\end{questao}
\begin{questao}{\nivel{N4} 09.43 — FGV}Uma plataforma oferece alternativa acessível somente mediante solicitação que demora 30 dias, enquanto usuários sem deficiência têm acesso imediato. A situação pode indicar:\begin{enumerate}[label=\Alph*)]\item barreira e tratamento desigual na fruição do serviço.\item plena equivalência de acesso.\item vantagem indevida do usuário com deficiência.\item ausência de qualquer dever de adaptação.\item impossibilidade de auditoria.\end{enumerate}\end{questao}
\begin{questao}{\nivel{N4} 09.44 — Cebraspe (C/E)}A avaliação de acessibilidade deve considerar não apenas existência formal de recursos, mas sua efetividade para uso autônomo e seguro.\end{questao}
\begin{questao}{\nivel{N4} 09.45 — FGV}Em avaliação de programa alinhado aos ODS, o auditor deve evitar:\begin{enumerate}[label=\Alph*)]\item verificar indicadores e resultados.\item tratar mera associação nominal a um ODS como prova suficiente de impacto.\item comparar metas e evidências.\item analisar grupos beneficiários.\item examinar sustentabilidade do resultado.\end{enumerate}\end{questao}
\begin{questao}{\nivel{N4} 09.46 — Cebraspe (C/E)}A Agenda 2030 pode ser usada como referência de planejamento e avaliação sem que isso signifique que cada uma de suas metas seja, isoladamente, título executivo judicial automático.\end{questao}
\begin{questao}{\nivel{N4} 09.47 — FGV}Um serviço concede prioridade aos doadores de sangue antes de gestantes e pessoas com deficiência. À luz da redação atual da Lei nº 10.048/2000, a regra está:\begin{enumerate}[label=\Alph*)]\item correta em qualquer situação.\item incompatível com a ordem especial prevista para doadores de sangue, cuja prioridade se dá após os demais beneficiários do rol.\item correta porque doador sempre possui prioridade absoluta.\item correta somente sem comprovante.\item fora do alcance da lei.\end{enumerate}\end{questao}
\begin{questao}{\nivel{N4} 09.48 — Cebraspe (C/E)}A ausência de barreira arquitetônica não exclui a possibilidade de discriminação por barreira atitudinal, comunicacional ou tecnológica.\end{questao}
\begin{questao}{\nivel{N4} 09.49 — FGV}Uma política racialmente neutra em sua redação produz exclusão sistematicamente maior de determinado grupo. Sob a ótica de igualdade material, o auditor deve:\begin{enumerate}[label=\Alph*)]\item ignorar o efeito porque não há menção explícita a raça.\item investigar dados, critérios e efeitos para verificar desigualdade e necessidade de correção.\item presumir que toda diferença é lícita.\item eliminar indicadores demográficos.\item avaliar somente custo financeiro.\end{enumerate}\end{questao}
\begin{questao}{\nivel{N4} 09.50 — Cebraspe (C/E)}Uma auditoria orientada por direitos humanos pode incorporar legalidade, acessibilidade, equidade e efetividade, examinando não apenas quanto foi gasto, mas quem conseguiu acessar o resultado da política.\end{questao}
